\chapter{Introduction}
\label{ch:intro}

% Seeded from PhD proposal Introduction (v10) with citations resolved against
% references.bib. Target 35-45 pp. Expand the background review (sec:intro_background)
% by drawing on the PDFs staged in proposal/writing/ch1_intro/.

\section{Motivation}
\label{sec:intro_motivation}

Additive manufacturing (AM) builds complex geometries directly from a digital model, enabling
rapid prototyping and distributed production through increasingly accessible
machines and materials~\cite{gibson2015additive,Beaman2020AdditiveMR,Yang2015AdditiveMD,yao2016cost,debroy2018additive,bahruni2018technology}.
It also gives designers the freedom to tailor each part to an individual
customer's requirements~\cite{ko2015design}.
Among powder-bed fusion techniques, selective laser sintering (SLS)
is widely adopted for polymer AM for its geometric flexibility and robust
mechanical performance~\cite{deckard1989method,jabri2023comprehensive,lupone2021process}.

SLS nonetheless presents considerable limitations. The process is
energy-inefficient: only a small fraction of the supplied thermal energy
directly contributes to part formation, while the majority heats unused
surrounding powder~\cite{jabri2023comprehensive,peyre2015sls,lupone2021process}.
Prolonged exposure of the powder to high temperature also degrades polymers such
as polyamide 12 (PA12), commonly known as Nylon~12, which can withstand only roughly
five to seven thermal cycles without powder refreshing before exhibiting reduced
mechanical strength, ductility, and
reliability~\cite{verbelen2016characterization,zhao2018crystallization,pham2008deterioration,choren2001sls,gornet2002characterization,zarringhalam2006effects}.
Maintaining a uniformly elevated build-chamber temperature further amplifies
energy losses while creating inconsistent thermal gradients that drive residual
stresses and warping~\cite{bourell2014performance,jabri2023comprehensive}. These
thermal side effects also lead to powder caking and semi-sintered boundaries,
increasing post-processing effort~\cite{xiao2024sls}.

Production throughput is similarly constrained. Build rate in SLS is limited
by the time required to expose every voxel with a scanned
laser~\cite{deckard1995improved}. Increasing speed requires either higher laser
power or multiple lasers, both of which introduce practical limitations. Excessive
energy density risks material degradation or ablation unless scan speed is raised,
but this is bounded by galvanometer performance. Adding lasers can increase
throughput, yet it also raises system cost and complexity
significantly~\cite{bourell2014performance,xiao2024sls}.

These limitations are not only technical; they carry significant economic
penalties. Machine depreciation, not raw material or tooling, dominates part cost
in SLS~\cite{lindemann2012lifecyclecosts,rickenbacher2013slmcostmodel}.
Capital equipment often exceeds
\$250{,}000,\footnote{Reported polymer powder-bed fusion system prices exceed
\$250{,}000; see \textcite{khajavi2014spareparts}.} yet builds are slow, producing
on the order of tens of parts per week for large, low-utilization builds, and
capacity utilization strongly governs the resulting per-part
cost~\cite{baumers2011capacityutilization}. To recover the capital
investment over a typical five-year lifespan, each part must carry a substantial
depreciation burden. Under this illustrative throughput, that burden works out to
more than \$100 per part on its own. Increasing build
rate is therefore not merely a matter of saving time, but of reducing part cost.
Faster, more efficient build processes are needed, and radio frequency additive manufacturing (RFAM) offers a
fundamentally different approach.

\section{The Radio-Frequency Additive Manufacturing (RFAM) Process}
\label{sec:intro_rfam}

RFAM volumetrically couples radio-frequency energy into doped regions of a
powder bed, selectively melting targeted areas without meaningfully raising the
temperature of the surrounding
material~\cite{allisonDissertation,allison2022volumetric,allison2022computational}.
RFAM uses a single additive, typically carbon black or graphite, in two roles at
once: it is the material patterned into the powder bed to define the part, and it is the
RF \emph{susceptor}, the phase that dissipates the applied field as heat. A binder-jet printhead deposits this susceptor
dopant layer by layer; the
resulting green part, the shaped but unsintered powder body, is then exposed to an radio frequency (RF) field for sintering. The
sintering mechanism is summarized in figure 1.1: the doped regions
couple to the field and dissipate energy through Joule heating, while the virgin
nylon remains nearly transparent at RF frequencies, so heat is generated only
within the patterned volume. The simulations in this work are exercised at the \SI{27.12}{\mega\hertz}
industrial, scientific, and medical (ISM) band, at which the available feedstock dielectric
data were measured; the hardware itself operates at \SI{13.56}{\mega\hertz}, whose
consequences for achievable resolution are weighed in chapter 7.

\begin{figure}[tb]
  \centering
  \includegraphics[width=\linewidth]{../build/figs/figures_ch1_rfam_process_overview.pdf.png}
  \caption{Figure 1.1. The RFAM sintering mechanism. Carbon/graphite-doped nylon powder is
  exposed to an RF field between parallel copper electrodes; the dopant
  selectively absorbs RF energy, melting the surrounding nylon and driving
  particle coalescence and part densification.}
  \label{fig:rfam_process}
\end{figure}

Decoupling energy input from surface scanning yields several advantages. Because
the entire doped cross-section is energized simultaneously rather than rastered,
throughput improves and thermal waste is reduced, and far fewer powder particles
are thermally cycled. Volumetric heating is also expected to improve material
homogeneity and reduce residual stress by eliminating the interlayer reheating
intrinsic to laser sintering~\cite{wudy2016filled}. The absence of heat in
non-targeted regions minimizes sintered artifacts in the surrounding powder and
reduces post-processing. Separately, the elimination of precision scanning optics allows
the architecture to scale without a corresponding rise in system complexity. The
underlying use of RF and microwave energy to heat lossy and doped materials is
well established outside of AM~\cite{wroe1998rfmicrowave,ferdous2017rfheating}.

Despite this potential, RFAM remains in its early stages, with open
questions around process control, achievable precision, and adaptability to
material systems beyond standard polymer/dopant combinations. This dissertation
addresses those questions by advancing RFAM process control, simulation, and
material extensibility.

\section{Background and Related Work}
\label{sec:intro_background}

\subsection{Polymer Powder-Bed Sintering}
\label{sec:intro_bg_sls}

Powder-bed fusion of polymers is dominated by selective laser sintering
(SLS), introduced by Deckard in the late 1980s~\cite{deckard1989method} and
since developed into a mature production
process~\cite{gibson2015additive,Beaman2020AdditiveMR,lupone2021process}. In
SLS, a scanned laser fuses successive powder layers held just below the
polymer melt temperature; the resulting parts have strong, isotropic mechanical
performance and good geometric flexibility~\cite{jabri2023comprehensive,peyre2015sls,xiao2024sls}.
The degree of particle melt achieved in the powder bed is the principal determinant
of final SLS part quality~\cite{Zarringhalam2009}. The
process is nonetheless constrained on three fronts. Energetically, only a small
fraction of supplied heat forms the part while the bulk maintains the elevated
chamber temperature, and the scanned exposure caps throughput~\cite{lupone2021process,deckard1995improved}.
Thermally, the elevated bed and repeated heating drive residual stress, warping,
and powder caking~\cite{bourell2014performance}.
Economically, slow builds on capital-intensive machines make depreciation the
dominant per-part cost~\cite{khajavi2014spareparts,baumers2011capacityutilization,lindemann2012lifecyclecosts,rickenbacher2013slmcostmodel,yao2016cost}.
These limitations (quantified in section 1.1) motivate the
volumetric, non-scanned alternative pursued in this dissertation.

\subsection{Thermal and Material Behavior of Polyamide-12}
\label{sec:intro_bg_pa12}

The processability of Nylon-12 (PA12), the workhorse SLS polymer and the
matrix used here, is governed by its semicrystalline thermal behavior. A narrow
window between melt-onset and recrystallization, characterized by differential
scanning calorimetry (differential scanning calorimetry (DSC)), defines the usable processing
range~\cite{zhao2018crystallization,affolter2001interlaboratory}, and the polymer
degrades after only a few thermal cycles, losing strength and ductility as powder
ages~\cite{verbelen2016characterization,zarringhalam2006effects,pham2008deterioration,choren2001sls,gornet2002characterization}.
Particle size and packing further govern layer uniformity and final
density~\cite{mokrane2018process}. These same DSC-measured transitions
(onset \SI{171}{\celsius}, peak \SI{180.8}{\celsius}, latent heat
\SI{101.7}{\kilo\joule\per\kilogram}, from the RFAM feedstock characterization of
Allison~\cite{allisonDissertation}) parameterize the apparent-heat-capacity melt
model used in the RFAM simulations of chapter 4, making PA12's
material thermodynamics a thread that runs through the entire work.

\subsection{Volumetric and Radio-Frequency Heating}
\label{sec:intro_bg_rf}

RFAM belongs to a broader shift away from serial, point-by-point energy delivery
toward parallel exposure of an entire region, in which the deposited energy field is
treated as a designed quantity. Tomographic volumetric additive manufacturing is the
clearest expression of that shift, solving for a set of projections before exposure so
that a whole volume solidifies at once~\cite{kelly2019volumetric}, and it now extends
to high-temperature inorganic parts through a subsequent furnace sinter step, a
conceptual parallel to RFAM's print-then-fuse sequence rather than the same
field-driven mechanism~\cite{toombs2022silica,madridwolff2023vamreview}.
Using radio-frequency and microwave energy to heat lossy and doped materials is
long established outside additive
manufacturing~\cite{wroe1998rfmicrowave,ferdous2017rfheating}, and within AM,
volumetric fusion of graphite-doped Nylon-12 under an RF field was demonstrated by
Allison and co-workers~\cite{allison2022volumetric,allisonDissertation}, whose
\SI{27.12}{\mega\hertz} material-property data the present simulations build on. The
governing physics, the prior RFAM results, and where recent microwave
volumetric-heating work sits relative to them are developed in chapter 2;
this dissertation advances that line by coupling a custom hardware platform
(section 3.1) with a predictive electromagnetic--thermal model and
simulation-guided compensation (chapter 5). RFAM draws on three research
lines, the volumetric-fusion paradigm~\cite{kelly2019volumetric}, concentration-tunable
RF and dielectric heating~\cite{vashisth2021rfsusceptors}, and coupled multiphysics
modeling of field-driven sintering~\cite{maniere2017coupled}.

\subsection{Binder Jetting and Dopant Deposition}
\label{sec:intro_bg_bjam}

In RFAM the susceptor dopant is patterned by binder jetting (selectively jetting
liquid onto a powder bed), a process introduced at MIT in the early 1990s and now
applied across metals, ceramics, and
polymers~\cite{mostafaei2021binderjetting,miyanaji2016porcelainBJ,du2020ceramicbj}.
Print fidelity is governed by droplet--powder interaction: a droplet impacts at
several meters per second, spreads, and infiltrates the pore network following
Washburn-type capillary flow, with the balance between spreading and infiltration
set by binder viscosity, surface tension, and powder
packing~\cite{Derby2010InkjetReview,colton2021droplet,holman2002spreading,barui2020inkpowder}.
High-speed X-ray imaging has resolved these dynamics in real
time~\cite{parab2019xray}. Small placement and volume errors accumulate into
banding, edge growth, and dimensional drift~\cite{Zhao2023BinderJetDefects,zago2024mbjaccuracy,Thalheim2023DropPlacement},
which is precisely what the metrology pipeline of section 3.3 is designed
to quantify.

\subsection{Process Compensation and Inverse Methods}
\label{sec:intro_bg_compensation}

Correcting process-induced non-uniformity is the core of chapter 5.
The primary designed strategy is \emph{functional grading} of dopant
concentration, which tunes local conductivity to flatten the heating field,
recasting the heuristic conductivity-tuning idea of prior RFAM
modeling~\cite{allison2022computational} as an explicit inverse of a verified
forward solver. Two further levers are characterized and
bounded rather than recommended for deployment. A second software lever, \emph{geometry pre-warp},
holds the dopant uniform and pre-distorts the part boundary by adjoint inverse
design; a transfer accounting reduces it to one narrow surviving role, a
chamber-keyed anisotropic shrinkage transform complementary to grading. On the
\emph{hardware} axis, turntable field averaging is the strongest lever in simulation
but is ruled out by the contact-coupling constraint, motivating a multi-stage
switchable-electrode system as the buildable path (chapter 5).

\subsection{Metrology for Additive Manufacturing}
\label{sec:intro_bg_metrology}

Quantitative, low-cost diagnostics underpin the characterization pipeline of
section 3.3. Flatbed scanners have been calibrated as traceable
two-coordinate measuring instruments~\cite{devicente2015scanner} and used for AM
layer-contour verification~\cite{blanco2020layercontour}, with print-quality
measurement standardized in the graphics industry~\cite{ISO24790}. Benchmark
test artifacts (the NIST artifact~\cite{moylan2014testartifact},
ISO/ASTM~52902~\cite{ISOASTM52902}, and the reviews of Rebaioli and
Fassi~\cite{rebaioli2017benchmark} and de Pastre et
al.~\cite{depastre2020testartefacts}) define dimensional-capability assessment,
while in-line approaches span fringe-projection
profilometry~\cite{zhang2016fringeprofilometry}, piezoelectric
self-sensing~\cite{wang2020piezoselfsensing}, and deep-learning image-based defect
detection~\cite{li2024deeplearningAM,kim2023imagefaultsurvey}. The present work
targets the on-substrate jetting outcome itself, before sintering, at commodity
cost (section 3.3.1).

\section{Research Questions and Dissertation Structure}
\label{sec:intro_rqs}

This work is organized around three tasks, each forming a chapter.

\begin{description}[leftmargin=0em,style=nextline]
  \item[Task~1: Design and characterization of a custom RFAM platform
  (sections 3.1 and 3.3).] Two questions follow. (i)~How can a modular
  RF additive manufacturing system integrate binder-jet dopant deposition,
  multi-axis motion, high-viscosity ink delivery, and volumetric RF heating?
  (ii)~How do low-cost, ex-situ diagnostics compare with specialized
  instrumentation for assessing jetting fidelity and dopant distribution?
  \item[Task~2: Toward ideal parts via computational and hardware compensation
  (chapters 4, 5 and 7).] Two questions follow.
  (i)~How can coupled electromagnetic--thermal models predict RF heating
  non-uniformity, and how can functionally graded doping be designed from those
  models as the primary compensation lever, with geometry pre-warp and turntable
  field averaging characterized and bounded rather than recommended for
  deployment? (ii)~What process-capability envelope and design-for-RFAM rules
  follow once those models are inverted (chapter 7)?
  \item[Task~3: Material and modality extensibility (chapter 8).]
  What dielectric and thermal criteria make a material RF-responsive, and how far
  can RFAM extend beyond powder beds, to bath-printed silicones, RF-foamed
  elastomers, and programmable deformation?
\end{description}

\section{Contributions}
\label{sec:intro_contributions}

At a high level, this dissertation contributes: (i) a custom, modular RFAM
platform and a low-cost scanner-based diagnostic pipeline (sections 3.1 and 3.3);
(ii) the high-frequency electrothermal additive thermal resolver (HEATR) coupled electromagnetic--thermal simulation framework and
simulation-guided compensation: functionally graded dopant maps that flatten the
RF heating field as the primary software lever. The contribution here is
\emph{computational and mechanistic}, and it is delineated carefully against prior
work. Song et al.~\cite{Song2025RFAMUniformity} demonstrated functional grading in
RFAM experimentally; the contribution here is not that experimental demonstration
but the inverse-design formulation that recasts the heuristic conductivity-tuning
idea of Allison's RFAM modeling~\cite{allison2022computational} as an explicit
inverse of a \emph{verified} forward solver, together with the mechanistic theory that
explains the method: the fixed-point ``glass-ceiling'' convergence analysis, the
integral-mode escape, and a method-selection framework that systematizes when grading
is the right lever. A retrospective \emph{validation} that the prior heuristic and
experimental work did not perform then calibrates that solver against Allison's own recovered
IR thermography (section 4.3.1). A second, distinct software
lever delivered here is \emph{geometry pre-warp} by adjoint-based inverse design:
holding the dopant uniform and pre-distorting the input geometry so that the
as-built part lands on the fixed nominal target. The boundary perturbations are
driven by gradients obtained from an adjoint through the coupled transient solver,
and verified end-to-end against finite differences. Its merit is that this gradient
route converges where prior \texttt{fmincon}-based attempts at the same problem did
not, and its one deployable role is a chamber-keyed anisotropic
densification-shrinkage transform complementary to grading; its limit, shared with FGM, is
the reentrant field-shadow case, where no
boundary pre-distortion (and no dopant grading) can back-fill heating energy into a
self-shadowed recess: a genuine negative result that bounds what compensation can
reach, not an unfinished one. On the \emph{hardware} axis,
turntable field averaging is the strongest lever in simulation but is ruled out by
the contact-coupling constraint, motivating the multi-stage switchable-electrode
path. A geometry-driven rule, validated across an 18-geometry method-selection
campaign, predicts which lever applies to a given part
(chapters 4 and 5); (iii) a
physics-derived capability envelope and a design-for-RFAM rule set, with the binding
minimum-section-thickness constraint demonstrated in-engine and the remaining limits
physics-derived, synthesizing the Task-2 results into what RFAM can and cannot build
(chapter 7); and (iv) a forward-looking reframing of
computationally programmed RF heating as a platform for spatially-graded material
transformation, with candidate applications (chapter 8).
Each contribution is stated in detail at the end of its chapter, after the
necessary terms and methods have been introduced.
